\section{Application}
 
The system will be attached to three standard econometric workloads chosen to span
Bayesian estimation, resampling inference, and simulation-based testing. In each case
the attachment is the same: one conditioned block seeds the workload's generator, the
transcript records that block with its provenance and health state, and a second
machine replays the transcript and re-runs the workload.
 
\begin{itemize}
  \item \textbf{Bayesian vector autoregression.} A BVAR with a Minnesota prior
        \citep{litterman1986} estimated by Gibbs sampling on a standard macroeconomic
        dataset, where posterior forecast bands depend on the sampler's random draws.
  \item \textbf{Bootstrap impulse-response intervals.} Bias-corrected bootstrap
        confidence intervals for impulse responses \citep{kilian1998}, where interval
        endpoints depend on the resampling sequence.
  \item \textbf{Monte Carlo evaluation of a forecast comparison test.} The empirical
        size of the Diebold--Mariano test \citep{diebold1995} under a known data-generating
        process, where rejection rates depend on the simulated error sequences.
\end{itemize}
 
Each workload demonstrates three things: that the provenance layer attaches without
changes to the estimation code (G4), that the run replays exactly on a second platform
(G3), and, as an illustration rather than a finding, how wide the band of results is
across independently seeded runs.